\chapter{Background}
\label{chap:background}

This chapter introduces the conceptual foundations for the thesis: the software supply-chain security context in which open-source dependencies are assessed, open-source sustainability and repository health metrics, the limitations of metric-based project assessment, and the machine-learning concepts relevant to tabular inactivity-risk scoring.

\section{Software Supply-Chain Security}
\label{sec:background-supply-chain-security}

Software supply-chain security reduces risks that arise when software artifacts, build processes, dependencies, and distribution channels are not sufficiently controlled or verifiable. Two frameworks place dependencies in this scope: the NIST Secure Software Development Framework treats the acquisition and maintenance of well-secured components as part of secure development practice \cite{souppaya2022ssdf}, and the SLSA framework defines integrity and provenance requirements for artifacts and build processes \cite{slsa2025spec} --- together motivating dependency monitoring beyond simple version updates.

Open-source dependencies are especially relevant because they are reused across many downstream projects. Open-source software is defined by licensing conditions that allow use, modification, and redistribution under criteria specified by the Open Source Initiative \cite{osiOpenSourceDefinition}; these freedoms enable broad reuse but also mean that downstream users depend on external maintainers and public infrastructure, so the maintenance activity of an upstream repository becomes relevant for downstream risk assessment. Supply-chain security thus provides the practical context for inactivity-risk scoring: the goal is not to determine whether a dependency is secure in general, but to estimate whether publicly observable repository signals indicate an elevated risk of future inactivity, supporting monitoring, replacement planning, or prioritization of manual review.

This context also explains why the unit of analysis is a single repository rather than a dependency graph: enumerating a project's dependencies is the task of established software-composition-analysis tooling, so the inventory is treated as an external input (Section~\ref{sec:impl-architecture}), and supply-chain security is the reason to look beyond version numbers to the maintenance trajectory of upstream repositories.

\section{Open-Source Sustainability and Project Health}
\label{sec:background-oss-health}

Open-source sustainability describes whether a project can continue to provide value and maintain its software over time. Project health is commonly assessed through observable indicators such as activity, responsiveness, contributor participation, and release behavior, for which the CHAOSS project provides standardized metrics models \cite{chaossStarterProjectHealth}.

Repository activity --- commits, releases, tags, issues, pull requests, and contributor activity --- is one category of observable health, collectable from public platforms and package metadata. Prior research uses revision activity and project attributes to study open-source project survival \cite{robinson2022survival}; in this thesis such signals form the basis for predicting whether a dependency repository becomes inactive within a 12-month horizon.

Security-practice indicators form a second category of observable signals. OpenSSF Scorecard assesses open-source projects using automated security-related checks \cite{zahan2022scorecard}. These describe repository configuration and security posture rather than future maintenance activity, and Section~\ref{sec:related-security-practice} sets out why this thesis keeps them as parallel triage context instead of feeding them into the inactivity model.

Project health is nonetheless broader than any single metric \cite{goggins2021transparent}, which is why inactivity-risk scoring is treated here as a limited operational signal for dependency triage rather than as a judgement of project sustainability.

These categories --- commit, contributor, issue and pull-request, release, age-and-popularity, and security-practice signals --- are the vocabulary from which the feature groups of this thesis are built. They are introduced here only as categories; the features themselves, with their exact definitions, are developed in Section~\ref{sec:method-feature-engineering}, where the design decisions behind them can be stated alongside them.

\section{Metric Pitfalls and Limitations}
\label{sec:background-metric-pitfalls}

Although repository metrics are useful, they can be misleading when interpreted without context \cite{yehudi2023contextfree}: a metric value can mean different things depending on the project's type, maturity, governance model, and maintenance style.

A low commit count may indicate reduced maintenance or merely a stable project that needs few changes; a high issue count may signal an active community or an unresolved backlog. This thesis therefore does not treat single metrics as proof of health or failure, but combines multiple signals and evaluates them empirically.

Responsiveness metrics require particular caution, because time-to-first-response can be affected by automated bot replies in pull-request workflows \cite{hasan2023firstresponse}. Such indicators should therefore be interpreted carefully and, where possible, separated into human and automated activity (Section~\ref{sec:related-responsiveness}).

Public repository data may also be incomplete: development may move platforms, releases may appear outside the observed repository, maintainers may use channels GitHub does not capture, and package-to-repository mappings may be ambiguous. The empirical results must therefore be read as evidence from publicly observable data, not a complete representation of all project activity.

These limitations shape the inactivity label: rather than project death or abandonment, inactivity is operationalized as the absence of selected observable activity signals within a defined horizon, which keeps the label measurable and reproducible while acknowledging it is only a proxy for elevated maintenance risk.

\section{Machine Learning Basics for Tabular Risk Scoring}
\label{sec:background-ml-basics}

The prediction task is binary classification: at observation time \(t\) the model estimates the probability that a repository becomes inactive in \( [t, t + 12\ \text{months}] \) from features derived from public metadata and activity. Because inactivity is closely related to recent activity, a recency baseline --- ranking repositories by the time since their last observed commit or release --- is a required reference point. It is worth being precise about what such a baseline is, because it is easy to mistake it for a backward-looking description. It is not: the recency baseline solves exactly the same forward-looking task as the model. Both read only the repository state at \(t\) and both are scored against the outcome in \( (t, t + 12\ \text{months}] \), which neither has seen. The baseline simply makes its forecast by extrapolating a single variable --- a repository that has been silent for a long time is predicted to stay silent --- where the learned model combines many. The comparison is therefore like-for-like, and it is a demanding one, since the label is itself derived from activity and a long silence is genuinely informative about future silence. This is why a learned model is worth reporting only if it improves on the trivial extrapolation: any skill it shows beyond that point is the part that could not be obtained by looking at the last commit date alone. Logistic regression serves as an interpretable probabilistic baseline, while tree-based models such as gradient boosting capture the non-linear feature interactions common in tabular repository metadata \cite{james2021islr,breiman2001randomforests,friedman2001greedy}.

Because a model of this kind can fail in two independent ways, it has to be judged on two independent questions, and the metrics used throughout this thesis correspond to them directly. The first question is \emph{discrimination}: does the model put the repositories that actually became inactive above the ones that did not? The second is \emph{calibration}: when the model outputs a probability of \(0.8\), do roughly \(80\%\) of such repositories in fact become inactive? A model can do well on one and badly on the other, which is why both are reported.

Discrimination is measured here by \textbf{ROC-AUC}, the area under the receiver operating characteristic curve. Its most useful reading is not geometric but probabilistic: ROC-AUC is the probability that a randomly chosen inactive repository receives a higher risk score than a randomly chosen active one. A value of \(0.5\) means the ordering is no better than a coin flip, and \(1.0\) means every inactive repository is ranked above every active one. The measure depends only on the \emph{order} of the scores, not on their absolute values, and it requires no decision threshold --- which makes it the right measure for a ranked triage view, where the analyst works down a list rather than applying a cut-off. Under strong class imbalance the area under the precision-recall curve is the more informative summary, because ROC-AUC can look flattering when the negative class dominates \cite{saito2015precision}.

Calibration is measured by the \textbf{Brier score} and the \textbf{expected calibration error (ECE)}. The Brier score is the mean squared difference between the predicted probability and the actual outcome (coded as \(1\) for inactive and \(0\) for active); lower is better, and the useful reference point is what a model with no skill at all would achieve --- predicting the base rate \(\bar{p}\) for every repository scores \(\bar{p}(1-\bar{p})\), so any value well below that indicates real information. The ECE takes a different view: it groups predictions into probability bins and averages how far each bin's predicted probability sits from the outcome frequency actually observed in it, so it reports miscalibration in the units of the score itself. Calibration matters here for a specific reason. A risk score that is only a ranking cannot support a statement such as ``this dependency has a one-in-four chance of going quiet within a year,'' and it is exactly that kind of statement --- not the rank --- that a team would act on when deciding whether to plan a replacement \cite{niculescu2005probabilities}. The formal definitions of all four metrics, together with the composite quality score used in the model comparison, are given in Appendix~\ref{sec:method-metric-definitions}.

All of these are computed on a \emph{held-out} slice: rows set aside before fitting and never used to train the model or to calibrate its probabilities, so that the reported numbers describe performance on data the model has not seen. Finally, the split into training and held-out data should be time-aware rather than random, so that information from later periods cannot influence model development for earlier ones; this matches the intended use of historical information to predict future inactivity.
